\bookchapter{第二章：整理你的数字房间——文件管理}

\chapterintro{手机用久了会乱，好习惯从管理文件开始。就像你的房间需要定期整理一样，你的数字空间也需要打理。一个井井有条的手机，不仅找东西方便，运行也更流畅。}

\booksection{一、三层存储体系：你的数字仓库}

\begin{enumerate}
  \item \textbf{机身存储}：核心区，速度快，系统和App都在这。这是你手机最宝贵的空间，要精打细算。
  \item \textbf{SD卡（如果支持）}：扩展衣柜，速度慢点，适合放电影、照片等大文件。注意：现在很多旗舰机不支持SD卡了。
  \item \textbf{云存储}：保险柜+同步网盘。防丢失，也能在多个设备间同步。iCloud、百度网盘、OneDrive都是常用选择。
\end{enumerate}

\booksection{二、认识“文件管理器”}

安卓一般有个 [文件] 或 [文件管理] App，iOS是自带的 [文件] App。它们是你查看手机所有文件的“资源管理器”。

\booksubsection{常见文件夹含义：}

\begin{itemize}
  \item \textbf{DCIM}：相机拍摄的照片和视频。
  \item \textbf{Pictures}：其他图片，如截图、保存的图片。
  \item \textbf{Download}：浏览器下载的文件。
  \item \textbf{Music}：音乐文件。
  \item \textbf{Movies}：视频文件。
  \item \textbf{Documents}：文档文件。
  \item \textbf{Android}：安卓系统文件夹，里面是各App的数据，一般不要手动删除。
\end{itemize}

\booksection{三、文件系统结构：深入了解}

安卓基于Linux，文件系统有一定的层级结构。虽然普通用户不需要深入了解，但知道一些基本概念有帮助。

\booksubsection{存储分区：}

\begin{itemize}
  \item \textbf{/system}：系统分区，存放操作系统文件。普通用户无权访问，需要root权限。
  \item \textbf{/data}：数据分区，存放用户数据和应用数据。这是你日常使用的空间。
  \item \textbf{/sdcard 或 /storage/emulated/0}：用户可访问的存储空间，你的照片、下载文件都在这里。
\end{itemize}

\booksubsection{App数据存储：}

\begin{itemize}
  \item \textbf{内部存储}：App的私有空间，其他App无法访问。卸载App时会被删除。
  \item \textbf{外部存储}：App在用户可访问区域存储的文件，如照片、下载内容。卸载App时不一定会被删除。
\end{itemize}

\booksection{四、数据备份与恢复：给你的数据上保险}

数据无价，备份是必须养成的习惯。

\booksubsection{备份策略：}

\begin{enumerate}
  \item \textbf{云备份}：自动、方便，但需要网络和云空间。iCloud、华为云、小米云等都支持。
  \item \textbf{本地备份}：用数据线连接电脑，用官方工具备份。速度快、空间大，但需要手动操作。
  \item \textbf{双重备份}：云备份+本地备份，最安全。
\end{enumerate}

\booksubsection{备份什么？}

\begin{itemize}
  \item \textbf{照片和视频}：最占空间，也最珍贵。
  \item \textbf{通讯录}：存在云账号里最方便。
  \item \textbf{微信聊天记录}：用微信自带的备份功能。
  \item \textbf{重要文档}：存在云盘或电脑上。
\end{itemize}

\booksection{五、文件安全：保护你的隐私}

\begin{itemize}
  \item \textbf{敏感文件加密}：很多手机有"安全文件夹"或"私密空间"功能，可以把敏感文件放进去，需要密码或指纹才能访问。
  \item \textbf{删除不等于销毁}：普通删除只是标记空间可覆盖，数据还在。要彻底删除，需要用专门的安全擦除工具，或者多次覆盖写入。
  \item \textbf{云端安全}：选择信誉好的云服务商，开启双重认证，定期检查登录设备。
\end{itemize}

\booksection{六、空间不够？来场大扫除}

\begin{itemize}
  \item \textbf{清理缓存}：App的临时记忆，清了能腾地方。在设置 > 存储 里可以清理。
  \item \textbf{查找大文件}：用文件管理器按大小排序，揪出占空间的"元凶"。
  \item \textbf{删除残留}：卸载App后，有时会留下垃圾文件，手动清掉。
  \item \textbf{清理相册}：删除重复照片、模糊照片、不需要的截图。
  \item \textbf{清理微信}：微信是空间大户，定期清理聊天记录和缓存很有必要。
\end{itemize}

\booksection{七、常见文件格式全解析}

手机里各种文件后缀名让人头大，搞懂它们能帮你更好地管理文件：

\begin{center}
\begin{tabularx}{\linewidth}{|X|X|X|X|}
\hline
\textbf{格式} & \textbf{类型} & \textbf{特性} & \textbf{推荐打开方式} \tabularnewline
\hline
JPG/JPEG & 图片 & 有损压缩，体积小，兼容性最好。手机拍照默认格式。 & 相册/任意图片查看器 \tabularnewline
PNG & 图片 & 无损压缩，支持透明背景。截图、图标常用。体积比JPG大。 & 相册/任意图片查看器 \tabularnewline
MP3 & 音频 & 有损压缩，体积小，兼容性极好。最通用的音频格式。 & 任意音乐播放器 \tabularnewline
FLAC & 音频 & 无损压缩，音质与CD一致，体积比MP3大3-5倍。发烧友首选。 & 音乐播放器/专业播放App \tabularnewline
WAV & 音频 & 无压缩原始音频，音质最好但体积巨大。专业录音常用。 & 音乐播放器 \tabularnewline
M4A & 音频 & AAC编码，比MP3同码率音质更好。苹果生态常用。 & Apple Music/音乐播放器 \tabularnewline
MP4 & 视频 & 最通用的视频封装格式，兼容性极好。手机录制的视频默认格式。 & 相册/任意视频播放器 \tabularnewline
MKV & 视频 & 开源封装格式，支持多音轨、多字幕。电影资源常用。 & VLC/nPlayer \tabularnewline
DOC/DOCX & 文档 & Word文档。DOCX是压缩格式，体积更小。DOC是旧版格式。 & WPS/Office \tabularnewline
PPT/PPTX & 演示 & PowerPoint演示文稿。PPTX是压缩格式。 & WPS/Office \tabularnewline
XLS/XLSX & 表格 & Excel电子表格。XLSX是压缩格式。 & WPS/Office \tabularnewline
PDF & 文档 & 跨平台格式，排版固定不可编辑。合同、论文、电子书常用。 & Adobe Reader/WPS \tabularnewline
ZIP & 压缩 & 最通用的压缩格式，跨平台兼容好。 & 系统自带/解压App \tabularnewline
RAR & 压缩 & 压缩率比ZIP高，但需要专用软件解压。 & 解压App \tabularnewline
ISO & 镜像 & 光盘镜像文件，完整复制光盘内容。系统安装包常用。 & 解压App/虚拟光驱 \tabularnewline
HTML & 网页 & 网页文件，用浏览器打开。保存的网页常用此格式。 & 浏览器 \tabularnewline
\hline
\end{tabularx}
\end{center}

\begin{tcbpractice}{练一练：整理"下载"文件夹}
打开文件管理器 > 进入“下载”文件夹。看看里面都有啥？把没用的删掉，有用的分门别类（比如新建“工作资料”、“表情包”文件夹放进去）。顺便看看存储空间，了解哪些类型文件占用最多。

\end{tcbpractice}

\booksection{八、图片与文件系统管理：让资产井井有条}

上一节我们学会了“认识文件夹”，但光认识还不够。本节要讲的是\textbf{如何主动建立一套可长期执行的管理规范}，让你的文件随时找得到、不丢失、不被偷窥。

\booksubsection{文件分类与命名规范}

混乱的根源是没有规则。建立规则其实很简单，关键在于\textbf{坚持}。

\begin{itemize}
  \item \textbf{"时间+事件"命名法}：文件夹名字里带时间戳和事件描述，方便日后回忆与检索。例如\texttt{202606\_旅行票据}、\texttt{2025年终总结}、\texttt{2024孩子生日}。这样既能看到时间线，也能一眼知道内容。
  \item \textbf{三级分类体系}：
  \begin{itemize}
    \item \textbf{一级}：按文件类型（文档 / 图片 / 视频 / 音频）
    \item \textbf{二级}：按使用场景（工作 / 生活 / 学习）
    \item \textbf{三级}：按具体事件或项目（如"2026Q2\_项目A"）
  \end{itemize}
  \item \textbf{月度清理计划}：每月固定一个日期（建议发工资那天或月初1号），对"下载"目录进行集中整理：\textbf{归档有用的、删除无用的、转存珍贵的}。一年 12 次小整理，比攒一年大扫除轻松得多。
\end{itemize}

\begin{tcbpractice}{本月任务清单}
挑一个空闲的周末，按以下顺序操作：

\begin{enumerate}
  \item 在文件管理器中创建\texttt{2026Q2\_项目A}形式的本月工作文件夹；
  \item 把所有"下载"目录里本周的文件过一遍，留下有用的、归档、删除没用的；
  \item 给今年的家庭照片建立\texttt{2025家庭}、\texttt{2026家庭}这种年度文件夹。
\end{enumerate}

\end{tcbpractice}

\booksubsection{冗余文件处理策略}

手机用久了会“长胖”，多半是冗余文件在作怪。

\begin{itemize}
  \item \textbf{72 小时自动清理}：微信、QQ 等社交 App 接收的临时文件（表情包、聊天图片、小视频），大多看一次就没用了。在 [微信] > [设置] > [通用] > [存储空间] 里开启"自动清理"，把过期缓存设成 72 小时。手机瞬间清爽不少，还不丢重要记录。
  \item \textbf{存储空间定向清理}：安卓 [设置] > [存储] > [清理加速]，iOS [设置] > [通用] > [iPhone 存储空间]，系统会按"缓存 / 临时文件 / 可卸载 App / 大文件"分类列出。重点清理：(1) 缓存文件；(2) 已过期安装包（升级 App 后留下的旧版安装包）；(3) 重复文件（同名同尺寸的照片）。
  \item \textbf{自动分类规则}：多数品牌相册 App 支持智能分类：人物、风景、文档、截图。开启后，新拍的合影会自动归入"人物 > 张三"，扫描的合同会自动归入"文档"。\textbf{配合每月整理，可极大降低"找不到东西"的概率}。
\end{itemize}

\booksubsection{图片安全处理规范}

图片是手机里最“敏感”的数据类型——一张随手拍就能暴露你的位置、设备、时间。我们建议\textbf{建立“双重保护”工作流}：

\begin{tcbwarn}{发送证件照/敏感截图的标准流程}
\begin{enumerate}
  \item \textbf{可见信息打码}：用系统自带或第三方工具（如美图秀秀、Snapseed）\textbf{涂抹身份证号中间几位、车牌号、门牌、姓名、金额}等关键字段；
  \item \textbf{元数据剥离}：发送前在相册中点"分享"或"导出"，系统会询问是否保留位置信息，选择\textbf{"不保留"}。这样拍摄设备型号、GPS 坐标、拍摄时间就不会跟着图片一起发出去；
  \item \textbf{临时水印}（可选）：在敏感截图上加"仅供XX办理使用"的文字水印，避免被滥用。
\end{enumerate}

\end{tcbwarn}

元数据（EXIF）清理的常用方法（基础操作，深度防护见第六章“照片脱敏处理”）：

\begin{itemize}
  \item \textbf{系统自带}：iOS [设置] > [隐私与安全性] > [相机]，可关闭"精确位置"；分享时选"仅当前照片"会\textbf{自动剥离大部分 EXIF}。
  \item \textbf{安卓}：多数品牌相册 App 在"分享"时会提示"是否移除位置信息"，勾选即可。
  \item \textbf{专业工具}：ExifTool（电脑端）、Metapho（iOS）、Photo Metadata Remover（安卓）可以彻底清除所有元数据。
\end{itemize}

对于需要\textbf{长期保存}的敏感图片（身份证、银行卡、户口本扫描件、护照），建议存入\textbf{加密相册}：

\begin{itemize}
  \item iPhone：[照片] > 选中照片 > 隐藏（仅 Face ID 验证后可查看）；
  \item 华为/小米：[相册] > 长按照片 > "移入保险柜"；
  \item 第三方加密 App：\textbf{360 安全相册}、\textbf{腾讯相册管家（隐私相册）}、\textbf{KeepSafe}等，单独密码保护，隐蔽性更好。
\end{itemize}
